\chapter*{Передмова автора}

\section*{Присвята}

У колофоні зазначено, що цей посібник присвячено всім державним службовцям України. Це означає, що як невід'ємна частина соціоінформаційної системи державного устрою вони є основними кінцевими користувачами нашої платформи. Саме для них — майбутніх розробників, операторів, архітекторів, керівників та аналітиків — і написано цей твір. Їхня щоденна праця забезпечує стабільність та розвиток державних інституцій, а впровадження новітніх технологічних рішень має на меті полегшити й оптимізувати цей надзвичайно важливий процес.

\section*{Ідея виникнення та історія розвитку системи}

Ідея розробки ERP/1 зародилася одразу після того, як я став на шлях підприємництва у 2005 році. Тоді я стояв перед вибором платформи для побудови моделі, котру планував впроваджувати, досліджуючи OCaml, Haskell і Erlang. Зрештою, моїм критеріям мотивації та технічним вимогам найкраще відповідала саме екосистема Erlang. Так розпочався мій професійний шлях: спершу був проєкт турецької соціальної мережі, згодом — робота у ПриватБанку, створення архітектури для месенджера NYNJA, навчання в аспірантурі з розробкою AXIO/1 і, нарешті, масштабні державні та міністерські системи.

Окремі частини системи, зокрема клієнтські компоненти для TWAIN-сканування та взаємодії з документами Word і Excel, було написано мовами C\# та F\#, оскільки мій найперший професійний досвід був тісно пов'язаний зі стеком .NET. Ось уже понад 18 років ERP/1 успішно обслуговує потреби замовників різного рівня, а фундаментальні ідеї платформи за цей час поширилися на екосистеми багатьох сучасних технологій, таких як LiveView, HTMX, Turbo, і навіть вплинули на такі старіші вебфреймворки, як Ocsigen, WebSharper, Lift та Nitrogen.

Сьогодні ERP/1 на базі Elixir повністю охоплює всі технологічні рівні державних інформаційних систем. При цьому вартість володіння платформою — що є одним із головних показників якості — зведено до абсолютного мінімуму. Це стало можливим завдяки тому, що всю систему написано єдиною мовою програмування і вона здатна повноцінно функціонувати на одному вузлі в межах однієї віртуальної машини. Для порівняння: зібраний код кожного з базових модулів ERP/1 легко вміщується на одній дискеті формату 3.5"\ ємністю 2.88 МБ.

Ефективність подібних підходів та низьку вартість володіння ними давно довели відомі продуктові компанії: WhatsApp, ПриватБанк, RabbitMQ, Basho, AdRoll, розробники Crytek Warface. Яскравим прикладом є компанія WhatsApp: на момент продажу гіганту Facebook її інфраструктуру обслуговували лише 6 інженерів, а вартість компанії сягала 19 мільярдів доларів. Системи на базі ERP/1 побудовано з використанням тієї ж самої надійної виробничої технології.


\section*{Розкриття компонентів основної мотивації}

Під час розробки системи ми послідовно застосовували філософію мінімалізму в галузі програмного забезпечення\footnote{\url{https://slides.com/maximsokhatsky/minimal/}}, про яку я вперше детально розповідав на одному з київських технологічних форумів ще 2013 року. Якщо викласти її суть стисло, вона зводиться до низки оптимізаційних критеріїв у вимірах часу, якості, людського та ресурсного капіталу:

1) зменшення часу виконання та розробки поряд зі значним подовженням життєвого циклу системи і підвищенням її ефективності. Основні метрики: Computation/Time, Weeks/Release, Feature/Hour, Hour/Lifetime;
2) подолання надмірної складності та зниження вартості розробки разом зі збільшенням надійності та концептуальної простоти. Основні метрики: Bugs/Code, Failures/Period, Cost/Support, Time/Fix;
3) усунення невизначеності на користь максимальної зрозумілості, відкритості та передбачуваності коду. Основні метрики: Committers/App, Issues/Feature, Msgs/Issue, Commits/User;
4) скорочення витрат на інфраструктурне обслуговування і додаткові ресурси задля значного збільшення обсягів обробки даних і обчислювальної спроможності. Основні метрики: Cost/Information, Machines/User, Data/User, Size/Features, Requests/Time, Information/User.

Якщо звести всі ці показники до єдиної всеосяжної максими, вона звучатиме так: ідеальний код — це той, котрий найлегше верифікувати (візуально, інструментально чи суворо математично), що, у свою чергу, неодмінно зумовлює його фундаментальну мінімалістичність. Наприклад, повна формальна верифікація рівня Г7 згідно з класифікатором ТЗІ (технічного захисту інформації) є недосяжною для таких систем, як типові SQL-бази даних. Адже практично неможливо математично довести всі без винятку властивості закритих пропрієтарних платформ на кшталт Oracle чи навіть потужних, але вкрай складних відкритих продуктів зразка PostgreSQL. Тому для підтвердження того, що система зберігає свої задані властивості, інженерам доводиться застосовувати безпосередній контроль типів і стовпців у таблицях, а всі функції згортки (рекурсори й фолди) власноруч писати в коді, наче це роблять курсорні генератори запитів у SQL. Тим не менш, до стандартного арсеналу Ericsson/OTP входить розподілена реляційна база даних Mnesia. І хоча вона теж не позбавлена певної академічної складності, її дизайн елегантно обмежений простим інтерфейсом, що гармонійно вписується в парадигму типізації System F.

\section*{Коментарі та настанови для слухачів курсу}

Цей навчальний курс є унікальним передусім тому, що предмет нашого дослідження — інформаційна система ERP/1 — відзначається продуманою формалізацією і чудово підходить як дидактичний матеріал для навчання фахівців. Разом із цим, ця ж система може похвалитися багатою історією успішних промислових впроваджень і безкомпромісною підтримкою сумісності.

Курс логічно поділено на 12 частин, кожна з яких послідовно розкриває певний етап розробки програмного забезпечення відповідно до стандартів ISO-9001. Матеріал охоплює повний цикл вітчизняного виробництва: від системного аналізу предметної області та роботи зі складною нормативно-юридичною документацією, через формування високорівневої архітектури, створення ескізного проєкту (з оглядкою на модель OSI та парадигму Закмана), — і аж до низькорівневих нюансів. До таких належать, зокрема, внутрішні механізми перевірки кваліфікованого електронного підпису (КЕП) та симетричного шифрування за алгоритмами ДСТУ-4145, реалізовані повністю автономно без залучення зовнішніх допоміжних криптографічних бібліотек. Завдяки високій щільності інформації та прикладній спрямованості цей посібник може стати вкрай корисним підручником і для фахівців напряму електронного урядування та цифрової трансформації державного апарату.

Матеріал створено з суворим урахуванням жорстких вимог, що висуваються до архітектури EDGE-офісів. Тут розглядається робота з системами, яким притаманні підвищені пороги допуску щодо розміру компільованої кодової бази, затримок мережевої взаємодії (latency), стійкості до відмов, рівня розподіленості, вартості подальшої підтримки, обчислювальної ефективності, багатства функціональності та експлуатаційної відповідності актуальним стандартам.

\section*{Подяки}

Я хотів би висловити найщирішу вдячність усім своїм вчителям, наставникам та академічним колегам, чия праця часто залишається недооціненою,
адже у глобальному масштабі інститут педагогіки переживає відчутну кризу поваги до професії викладача.
Особлива подяка викладачам математики, інженерії та філософії, які свого часу прищепили мені фундаментальне прагнення до логічної
чистоти та краси архітектурних рішень. Крім того, неабияк дякую всім контриб'юторам платформи ERP/1,
які невтомно вдосконалювали цей продукт, а також своїм рідним і близьким за їхню багаторічну безмежну підтримку.

\chapter{Вступ}

Цей посібник докладно описує формальну модель та програмну архітектуру сучасного функціонального верифікованого мовного забезпечення. Його призначення — слугувати фундаментом для побудови складних інфраструктурних процесінгових систем, орієнтованих передусім на державну модель управління процесами. Платформа створена з прицілом на можливість прозорого проведення глибокого зовнішнього аудиту; вона сумісна з різноманітними видами розподілених сховищ, сучасними телекомунікаційними та реєстровими фреймворками, базовими інтернет-утворюючими сервісами, і головне — ідеально підлаштована під завдання автоматизації захищених автономних офісів та великих державних підприємств України у цілому.

Система управління державними підприємствами ERP/1, що презентується в цьому посібнику, є не тільки ідіоматичним зразком побудови сучасних інформаційних державних та комерційних екосистем. Вона насамперед визначає сувору формальну специфікацію та демонструє її надійну імплементацію (що налічує вже безліч реальних національних впроваджень) для потреб інноваційних, високооптимізованих підприємств, які вимагають максимально точних інструментів контролю операцій і беззаперечних гарантій цілісності даних.

Телекомунікаційна платформа Erlang/OTP, створена в надрах компанії Ericsson, успішно та безперебійно застосовується у світовій індустрії передовими мобільними операторами вже понад 30 років, а її віртуальна машина і досі заслужено вважається однією з найнадійніших у галузі. Інформаційні системи управління підприємствами на її базі вже багато років використовуються не лише в телекомі, але й в елітарній банківській сфері, високошвидкісному фінансовому процесінгу транзакцій, стійких розподілених брокерах повідомлень та в IoT-секторі. Ви в будь-який час можете переглянути роботу наших демо-модулів системи ERP/1, розгорнутих у спеціальному захищеному середовищі з власним автономним центром генерації та випуску ECC X.509 сертифікатів. На сторінках цієї книги ви знайдете вичерпний перелік ключових модулів системи та опис фундаментальних сутностей схеми.

В основі нашої концепції лежить ідея створення універсальної базової платформи для розробки та стабільного функціонування різноманітних інформаційних реєстрів і розподілених баз (чи банків) даних будь-яких масштабів: від невеликих міжсистемних довідників і спеціалізованих класифікаторів до високонавантажених корпоративних, місцевих та масштабних державних інформаційних ресурсів. Переконаний, що цей посібник буде вкрай корисним для всіх, хто бажає глибоко зрозуміти принципи роботи, архітектуру та методики створення надійних інформаційних платформ, які сьогодні активно застосовуються як у вітчизняному державному, так і в комерційному секторах.

\section*{Список скорочень}

\begin{description}
    \item[ВККС] — Вища кваліфікаційна комісія суддів України.
    \item[ВРП] — Вища рада правосуддя.
    \item[ГРД] — Громадська рада доброчесності.
    \item[ГРМЕ] — Громадська рада міжнародних експертів.
    \item[ГСЦ] — Головний сервісний центр МВС.
    \item[ДМС] — Державна міграційна служба України.
    \item[ДПС] — Державна прикордонна служба України.
    \item[ДСА] — Державна судова адміністрація.
    \item[ДСНС] — Державна служба України з надзвичайних ситуацій.
    \item[ЕСОЗ] — Електронна система охорони здоров'я МОЗ.
    \item[ЄІС] — Єдина інформаційна система МВС.
    \item[ЄРЗ] — Єдиний реєстр зброї НПУ.
    \item[ЄСІКС] — Єдина судова інформаційно-комунікаційна система.
    \item[КЕП] — Кваліфікований електронний підпис.
    \item[КЗІ] — Криптографічний захист інформації.
    \item[КМУ] — Кабінет Міністрів України.
    \item[КСЗІ] — Комплексна система захисту інформації.
    \item[МВС] — Міністерство внутрішніх справ України.
    \item[МОЗ] — Міністерство охорони здоров'я України.
    \item[МОН] — Міністерство освіти і науки України.
    \item[МОУ] — Міністерство оборони України.
    \item[НАН] — Національна академія наук України.
    \item[НГУ] — Національна гвардія України.
    \item[НДІ] — Науково-дослідний інститут.
    \item[НПА] — Нормативно-правовий акт.
    \item[НПУ] — Національна поліція України.
    \item[ОВВ] — Органи виконавчої влади.
    \item[ПФУ] — Пенсійний фонд України.
    \item[ССО] — Служба судової охорони / Сили спеціальних операцій.
    \item[СУСЗЦЗ] — Система управління силами та засобами цивільного захисту ДСНС.
    \item[ТЗІ] — Технічний захист інформації.
    \item[СТЗІ] — Система технічного захисту інформації.
    \item[ФП МТРЗ] — Функціональна підсистема матеріально-технічного та ресурсного забезпечення МВС.
    \item[ЦГЗ] — Центр громадського здоров'я.
    \item[ЦОВВ] — Центральний орган виконавчої влади.
    \item[ЦОД] — Центр обробки даних.
\end{description}
